\documentclass[aspectratio=169]{beamer}

\usetheme{Madrid}
\usecolortheme{default}
\setbeamertemplate{navigation symbols}{}

\usepackage{graphicx}
\usepackage{tikz}
\usepackage{ragged2e}

\definecolor{UGBlue}{RGB}{21,55,94}
\definecolor{UGLight}{RGB}{243,246,251}

\setbeamercolor{frametitle}{bg=UGBlue,fg=white}
\setbeamercolor{title}{fg=UGBlue}
\setbeamercolor{structure}{fg=UGBlue}

\title{Resource-Constrained Evaluation of Lightweight Concept Drift Detectors in Machine Learning}
\subtitle{An Empirical Evaluation of Lightweight Drift Detectors Under Resource Constraints}
\author{Stephen Kwaku Azongo}
\institute{Department of Computer Science, University of Ghana}
\date{Defense Presentation}

\newcommand{\templatecover}[5]{%
  \begin{frame}[plain]
    \begin{tikzpicture}[remember picture,overlay]
      \node[anchor=south west, inner sep=0] at (current page.south west)
      {\includegraphics[width=\paperwidth,height=\paperheight]{assets/ug_template_cover_bg.jpg}};

      \node[anchor=north, text=white, font=\bfseries\large] at ([yshift=-0.45cm]current page.north)
      {FINAL YEAR PROJECT PRESENTATION \& DEMO};

      \node[anchor=north, text=white, font=\bfseries\large, align=center, text width=0.80\paperwidth]
      at ([yshift=-1.35cm]current page.north) {#1};

      \node[anchor=north, text=white, font=\bfseries\normalsize, align=center]
      at ([yshift=-2.65cm]current page.north) {#2\\#3\\#4};

      \node[anchor=north, text=white, font=\bfseries\scriptsize, align=center]
      at ([yshift=-4.35cm]current page.north) {SUPERVISOR: #5\\8\textsuperscript{TH} OCTOBER 2026};
    \end{tikzpicture}
  \end{frame}
}

\begin{document}

\templatecover
{Resource-Constrained Evaluation of Lightweight Concept Drift Detectors in Machine Learning}
{Stephen Kwaku Azongo}
{ID: 11027858}
{Department of Computer Science, University of Ghana}
{Prof. KOFI SARPONG ADU-MANU}

\begin{frame}{OUTLINE OF PRESENTATION}
  \begin{enumerate}
    \item Introduction
    \item Aims and Objectives
    \item Related Work
    \item Proposed Solution (Methodology)
    \item Key Results
    \item Conclusion
  \end{enumerate}
\end{frame}

\begin{frame}{Introduction}
  \justifying
  \textbf{Problem context:} Machine learning models in production face concept drift, where the joint distribution \(P(X,y)\) changes over time and degrades prediction quality.
  \vspace{0.4em}

  \textbf{Research gap:}
  \begin{itemize}
    \item Existing drift benchmarks usually hold deployment resources constant.
    \item Real systems must trade off monitoring frequency, detection window size, and reference memory.
    \item It remains unclear which constraint matters most in practical detection performance.
  \end{itemize}

  \textbf{Study focus:} Evaluate lightweight unsupervised detectors under realistic resource constraints across financial, telecom, and industrial datasets.
\end{frame}

\begin{frame}{Aims and Objectives}
  \textbf{Aim:}
  \begin{itemize}
    \item To determine whether monitoring configuration is more consequential than detector algorithm under resource-constrained deployment.
  \end{itemize}

  \textbf{Specific Objectives:}
  \begin{enumerate}
    \item Compare KS, PSI, and Wasserstein detectors over 8,100 experimental runs.
    \item Quantify the effect of monitoring frequency, window size, and reference size on delay and missed detection.
    \item Evaluate sensitivity across sudden, gradual, and recurring drift.
    \item Assess downstream model robustness (Random Forest, Logistic Regression, LightGBM) under controlled drift.
  \end{enumerate}
\end{frame}

\begin{frame}{Related Work}
  \begin{itemize}
    \item \textbf{Concept drift foundations:} Gama et al.\ formalize drift types and adaptive monitoring needs.
    \item \textbf{Empirical detector benchmarks:} Barros \& Santos; Lukats et al.\ compare many detectors but do not model deployment constraints as primary factors.
    \item \textbf{Distribution-based monitoring:} KS, PSI, and Wasserstein are practical, label-free, and lightweight for production settings.
    \item \textbf{Open gap addressed in this project:}
    \begin{itemize}
      \item No prior study systematically ranks resource constraints by their quantitative impact on detection outcomes.
    \end{itemize}
  \end{itemize}
\end{frame}

\begin{frame}{Proposed Solution (Methodology)}
  \begin{itemize}
    \item \textbf{Datasets:} UCI Credit Default, IBM Telco Churn, and UCI AI4I Predictive Maintenance.
    \item \textbf{Drift simulation:} Sudden, gradual, and recurring drift injected into stream portions after a 60/40 reference-stream split.
    \item \textbf{Detectors:} Kolmogorov-Smirnov test, Population Stability Index, and Wasserstein distance.
    \item \textbf{Constraint grid:}
    \begin{itemize}
      \item Monitoring frequency: \(\{10, 50, 100, 500, 1000\}\)
      \item Detection window size: \(\{50, 100, 500, 1000\}\)
      \item Reference size: \(\{100, 500, \text{unlimited}\}\)
    \end{itemize}
    \item \textbf{Analysis:} Kruskal-Wallis, pairwise Mann-Whitney U (Bonferroni), Spearman correlation, and effect sizes.
  \end{itemize}
\end{frame}

\begin{frame}{Key Results}
  \begin{itemize}
    \item Constraint configuration produced a \textbf{37-fold} range in detection delay, much larger than detector-to-detector differences.
    \item Gradual drift was significantly harder to detect than sudden/recurring drift (\(r \approx 0.30\), medium effect).
    \item Detection window size had the strongest individual impact; monitoring frequency was second.
    \item A 500-sample reference window performed comparably to unlimited reference in most settings.
    \item For post-drift model robustness, Random Forest degraded less than LightGBM (\(r=0.383\)) and Logistic Regression (\(r=0.185\)).
  \end{itemize}
\end{frame}

\begin{frame}{Conclusion}
  \begin{block}{Main Conclusion}
    In resource-constrained monitoring, \textbf{configuration matters more than algorithm} for practical drift detection performance.
  \end{block}

  \textbf{Deployment priority hierarchy:}
  \begin{enumerate}
    \item Configure monitoring regime first (frequency, window, reference).
    \item Account for drift type behavior, especially gradual drift.
    \item Then optimize detector and predictive model selection.
  \end{enumerate}

  \vspace{0.5em}
  \centering
  \large \textbf{Thank You}
\end{frame}

\end{document}
